% =============================================================================
% Capítulo 7: Anexos
% =============================================================================
\section{Anexos}
\label{sec:anexos}

\subsection{Anexo A: Diccionario de Datos}
\label{sec:anexo_diccionario}

Este anexo describe cada columna del esquema \bd{dwh}: tipo de dato, si admite
nulos, las llaves y reglas en que participa (PK: llave primaria; FK: llave
foránea; UQ: valor único; CK: regla \bd{CHECK}) y su descripción. Las
descripciones son los comentarios \bd{COMMENT ON} del propio script, y el
anexo se genera desde el catálogo de PostgreSQL con
\bd{herramientas/generar_diccionario.py}, por lo que coincide con la base
creada.

% =============================================================================
% Anexo A: Diccionario de Datos.
% Generado por herramientas/generar_diccionario.py desde el catálogo de
% PostgreSQL. No editar a mano: volver a ejecutar el script.
% =============================================================================

\subsubsection{Tabla \bd{fact_venta}}

Hechos de venta: una fila por transacción de compra de un cliente. \textbf{Origen:} Transacciones de las seis fuentes públicas de la Sección~\ref{sec:fuentes} (una fila por pasaje, viaje, envío o pedido), asignadas a clientes sintéticos.

\begin{table}[H]
  \centering
  \caption{Diccionario de datos de \bd{fact_venta}}
  \small
  \begin{tabularx}{\linewidth}{@{}>{\raggedright\arraybackslash}p{4.0cm}lcl>{\raggedright\arraybackslash}X@{}}
    \toprule
    \textbf{Columna} & \textbf{Tipo} & \textbf{Nulo} & \textbf{Clave} & \textbf{Descripción} \\
    \midrule
    \bd{id_venta} & BIGINT & No & PK & PK. Llave sustituta \\
    \bd{id_tiempo} & INTEGER & No & FK & FK a dim\_tiempo (fecha de la compra) \\
    \bd{id_cliente} & INTEGER & No & FK & FK a dim\_cliente \\
    \bd{id_servicio} & INTEGER & No & FK & FK a dim\_servicio \\
    \bd{id_zona_origen} & INTEGER & No & FK & FK a dim\_zona, rol origen \\
    \bd{id_zona_destino} & INTEGER & No & FK & FK a dim\_zona, rol destino \\
    \bd{id_canal_venta} & INTEGER & No & FK & FK a dim\_canal\_venta \\
    \bd{id_medio_pago} & INTEGER & No & FK & FK a dim\_medio\_pago \\
    \bd{numero_transaccion} & VARCHAR(30) & No & UQ & Dimensión degenerada: id del sistema de origen, prefijado por línea \\
    \bd{edad_cliente} & SMALLINT & Sí & CK & Edad al momento de la compra (no aditiva; NULL para empresas) \\
    \bd{cantidad} & INTEGER & No & CK & Medida aditiva: pasajes, bultos o unidades \\
    \bd{monto_bruto} & NUMERIC(12,2) & No & CK & Medida aditiva: monto antes de descuento (CLP) \\
    \bd{monto_descuento} & NUMERIC(12,2) & No & CK & Medida aditiva: descuento aplicado (CLP) \\
    \bd{monto_neto} & NUMERIC(12,2) & No & CK & Medida aditiva: monto pagado (CLP) = bruto - descuento \\
    \bd{distancia_km} & NUMERIC(10,2) & Sí & CK & Medida aditiva: distancia recorrida en km \\
    \bd{duracion_minutos} & INTEGER & Sí & CK & Medida aditiva: duración del servicio en minutos \\
    \bd{peso_kg} & NUMERIC(10,2) & Sí & CK & Medida aditiva: peso en kg (solo logística y delivery) \\
    \bottomrule
  \end{tabularx}
\end{table}

\subsubsection{Tabla \bd{dim_tiempo}}

Dimensión tiempo con granularidad diaria. \textbf{Origen:} Se genera con un script SQL (\bd{generate_series}) para 2023 a 2025; no requiere fuente externa.

\begin{table}[H]
  \centering
  \caption{Diccionario de datos de \bd{dim_tiempo}}
  \small
  \begin{tabularx}{\linewidth}{@{}>{\raggedright\arraybackslash}p{4.0cm}lcl>{\raggedright\arraybackslash}X@{}}
    \toprule
    \textbf{Columna} & \textbf{Tipo} & \textbf{Nulo} & \textbf{Clave} & \textbf{Descripción} \\
    \midrule
    \bd{id_tiempo} & INTEGER & No & PK, CK & PK. Fecha en formato AAAAMMDD (ej. 20250314) \\
    \bd{fecha} & DATE & No & UQ & Fecha calendario \\
    \bd{dia} & SMALLINT & No & CK & Día del mes (1-31) \\
    \bd{dia_semana} & SMALLINT & No & CK & 1 = lunes … 7 = domingo (ISO) \\
    \bd{nombre_dia} & VARCHAR(10) & No &  & Lunes … Domingo \\
    \bd{es_fin_de_semana} & BOOLEAN & No & CK & TRUE si es sábado o domingo \\
    \bd{semana_anio} & SMALLINT & No & CK & Semana ISO del año (1-53) \\
    \bd{mes} & SMALLINT & No & CK & Mes (1-12) \\
    \bd{nombre_mes} & VARCHAR(10) & No &  & Enero … Diciembre \\
    \bd{trimestre} & SMALLINT & No & CK & Trimestre (1-4) \\
    \bd{semestre} & SMALLINT & No & CK & Semestre (1-2) \\
    \bd{anio} & SMALLINT & No & CK & Año \\
    \bottomrule
  \end{tabularx}
\end{table}

\subsubsection{Tabla \bd{dim_cliente}}

Dimensión cliente unificada entre las seis líneas de negocio. \textbf{Origen:} Sintética: clientes residentes en Chile, con la distribución por comuna, sexo y edad del Censo de Población y Vivienda 2024.

\begin{table}[H]
  \centering
  \caption{Diccionario de datos de \bd{dim_cliente}}
  \small
  \begin{tabularx}{\linewidth}{@{}>{\raggedright\arraybackslash}p{4.0cm}lcl>{\raggedright\arraybackslash}X@{}}
    \toprule
    \textbf{Columna} & \textbf{Tipo} & \textbf{Nulo} & \textbf{Clave} & \textbf{Descripción} \\
    \midrule
    \bd{id_cliente} & INTEGER & No & PK & PK. Llave sustituta \\
    \bd{codigo_cliente} & VARCHAR(20) & No & UQ & Llave natural (RUT o documento, unificado entre líneas) \\
    \bd{nombre_cliente} & VARCHAR(100) & No &  & Nombre o razón social \\
    \bd{genero} & VARCHAR(20) & Sí &  & Género declarado (NULL para empresas) \\
    \bd{fecha_nacimiento} & DATE & Sí &  & Fecha de nacimiento (NULL para empresas) \\
    \bd{tipo_cliente} & VARCHAR(20) & No & CK & persona o empresa \\
    \bd{comuna_residencia} & VARCHAR(60) & Sí &  & Comuna de residencia \\
    \bd{region_residencia} & VARCHAR(60) & Sí &  & Región de residencia \\
    \bd{pais_residencia} & VARCHAR(60) & No &  & País de residencia \\
    \bd{fecha_registro} & DATE & Sí &  & Fecha de alta como cliente \\
    \bottomrule
  \end{tabularx}
\end{table}

\subsubsection{Tabla \bd{dim_servicio}}

Dimensión servicio: línea de negocio y medio de transporte. \textbf{Origen:} Catálogo de los servicios de cada línea de negocio, con los atributos que informa cada fuente (tipo y clase de servicio, medio de transporte).

\begin{table}[H]
  \centering
  \caption{Diccionario de datos de \bd{dim_servicio}}
  \small
  \begin{tabularx}{\linewidth}{@{}>{\raggedright\arraybackslash}p{4.0cm}lcl>{\raggedright\arraybackslash}X@{}}
    \toprule
    \textbf{Columna} & \textbf{Tipo} & \textbf{Nulo} & \textbf{Clave} & \textbf{Descripción} \\
    \midrule
    \bd{id_servicio} & INTEGER & No & PK & PK. Llave sustituta \\
    \bd{codigo_servicio} & VARCHAR(20) & No & UQ & Llave natural \\
    \bd{linea_negocio} & VARCHAR(50) & No & CK & Una de las 6 líneas de negocio del enunciado (nombre literal) \\
    \bd{medio_transporte} & VARCHAR(30) & No &  & bus, tren, taxi, camión, avión, moto/bicicleta \\
    \bd{tipo_servicio} & VARCHAR(20) & No & CK & pasajeros o carga \\
    \bd{clase_servicio} & VARCHAR(30) & No &  & Ej. económica, ejecutiva, express, estándar \\
    \bottomrule
  \end{tabularx}
\end{table}

\subsubsection{Tabla \bd{dim_zona}}

Dimensión zona geográfica, de rol doble (origen y destino). \textbf{Origen:} Tablas de zonas de cada fuente (zonas de taxi de Nueva York, puntos de ANTT, estaciones del Reino Unido, aeropuertos de EE.UU.) y comunas de Chile del Censo 2024.

\begin{table}[H]
  \centering
  \caption{Diccionario de datos de \bd{dim_zona}}
  \small
  \begin{tabularx}{\linewidth}{@{}>{\raggedright\arraybackslash}p{4.0cm}lcl>{\raggedright\arraybackslash}X@{}}
    \toprule
    \textbf{Columna} & \textbf{Tipo} & \textbf{Nulo} & \textbf{Clave} & \textbf{Descripción} \\
    \midrule
    \bd{id_zona} & INTEGER & No & PK & PK. Llave sustituta \\
    \bd{codigo_zona} & VARCHAR(20) & No & UQ & Llave natural \\
    \bd{comuna} & VARCHAR(60) & No &  & Comuna (No aplica en destinos internacionales) \\
    \bd{ciudad} & VARCHAR(60) & No &  & Ciudad \\
    \bd{region} & VARCHAR(60) & No &  & Región \\
    \bd{pais} & VARCHAR(60) & No &  & País \\
    \bd{continente} & VARCHAR(30) & No &  & Continente \\
    \bottomrule
  \end{tabularx}
\end{table}

\subsubsection{Tabla \bd{dim_canal_venta}}

Dimensión canal de venta. \textbf{Origen:} Canales de compra informados por las fuentes, normalizados a un catálogo común.

\begin{table}[H]
  \centering
  \caption{Diccionario de datos de \bd{dim_canal_venta}}
  \small
  \begin{tabularx}{\linewidth}{@{}>{\raggedright\arraybackslash}p{4.0cm}lcl>{\raggedright\arraybackslash}X@{}}
    \toprule
    \textbf{Columna} & \textbf{Tipo} & \textbf{Nulo} & \textbf{Clave} & \textbf{Descripción} \\
    \midrule
    \bd{id_canal_venta} & INTEGER & No & PK & PK. Llave sustituta \\
    \bd{codigo_canal_venta} & VARCHAR(20) & No & UQ & Llave natural \\
    \bd{nombre_canal_venta} & VARCHAR(40) & No &  & app, web, boletería, call center, convenio empresa \\
    \bottomrule
  \end{tabularx}
\end{table}

\subsubsection{Tabla \bd{dim_medio_pago}}

Dimensión medio de pago. \textbf{Origen:} Tipos de pago informados por las fuentes, normalizados a un catálogo común.

\begin{table}[H]
  \centering
  \caption{Diccionario de datos de \bd{dim_medio_pago}}
  \small
  \begin{tabularx}{\linewidth}{@{}>{\raggedright\arraybackslash}p{4.0cm}lcl>{\raggedright\arraybackslash}X@{}}
    \toprule
    \textbf{Columna} & \textbf{Tipo} & \textbf{Nulo} & \textbf{Clave} & \textbf{Descripción} \\
    \midrule
    \bd{id_medio_pago} & INTEGER & No & PK & PK. Llave sustituta \\
    \bd{codigo_medio_pago} & VARCHAR(20) & No & UQ & Llave natural \\
    \bd{nombre_medio_pago} & VARCHAR(40) & No &  & efectivo, débito, crédito, transferencia, billetera digital \\
    \bottomrule
  \end{tabularx}
\end{table}


\subsection{Anexo B: Reproducción del Entorno}
\label{sec:anexo_entorno}

\subsubsection{Ejecución de los Scripts en Cualquier Equipo}

Los scripts \bd{01_ddl.sql} y \bd{99_validacion.sql} son autónomos: no leen ni
escriben archivos, no dependen de otros archivos y se ejecutan desde cualquier
carpeta, en Windows, Linux o macOS. Solo requieren PostgreSQL 17 y una base de
datos existente; se recomienda una base nueva, porque \bd{01_ddl.sql} borra y
vuelve a crear el esquema \bd{dwh}.

\begin{enumerate}[leftmargin=*]
  \item Crear una base de datos UTF-8, por ejemplo \bd{rapiditos}.
  \item En pgAdmin, abrir \bd{01_ddl.sql} en el Query Tool de esa base y
        ejecutarlo completo con F5. Con \bd{psql}, desde la carpeta en que
        estén los scripts:
\begin{lstlisting}[style=consola]
psql -U <usuario> -d rapiditos -v ON_ERROR_STOP=1 -f 01_ddl.sql
psql -U <usuario> -d rapiditos -P pager=off -f 99_validacion.sql
\end{lstlisting}
  \item Comparar el resultado de \bd{99_validacion.sql} con la
        Tabla~\ref{tab:verificacion}: la última consulta debe devolver cero
        filas.
\end{enumerate}

Para Windows se entrega además la guía \bd{guia_ejecucion_windows.pdf}, que
detalla paso a paso la instalación de PostgreSQL y pgAdmin, la ejecución con
\bd{psql} y una alternativa con Docker Desktop que no requiere instalar
PostgreSQL ni ningún otro archivo del proyecto.

\subsubsection{Entorno de Desarrollo del Grupo}

El grupo trabaja con el repositorio completo del proyecto, que define
PostgreSQL y pgAdmin en \bd{docker-compose.yml} e incluye las pruebas
automáticas. Estos pasos requieren ese repositorio y se ejecutan desde su
raíz, igual en Linux y en Windows; solo cambian la copia del archivo
\bd{.env} y la activación del entorno virtual de Python, indicadas en los
comentarios.

\begin{enumerate}[leftmargin=*]
  \item Preparar las credenciales y levantar PostgreSQL y pgAdmin. En
        \bd{.env} se definen las contraseñas de ambos servicios.
\begin{lstlisting}[style=consola]
cp .env.example .env      # Windows: copy .env.example .env
docker compose up -d
docker compose ps         # ambos servicios deben aparecer activos
\end{lstlisting}
  \item Crear el modelo multidimensional; pgAdmin queda en
        \bd{http://localhost:5050}.
\begin{lstlisting}[style=consola]
docker compose exec postgres psql -U rapiditos -d rapiditos \
    -v ON_ERROR_STOP=1 -f /sql/01_ddl.sql
\end{lstlisting}
  \item Crear el entorno de Python y ejecutar las pruebas automáticas, que
        usan una base de prueba aparte y nunca modifican \bd{rapiditos}.
\begin{lstlisting}[style=consola]
python -m venv .venv
source .venv/bin/activate # Windows: .venv\Scripts\activate
pip install -r requirements.txt
pytest -v
\end{lstlisting}
  \item Opcionalmente, regenerar el diccionario de datos del Anexo~A y las
        capturas del informe.
\begin{lstlisting}[style=consola]
python herramientas/generar_diccionario.py
python -m playwright install chromium   # una sola vez
python herramientas/capturar_pgadmin.py
\end{lstlisting}
\end{enumerate}

\subsection{Anexo C: Repositorio del Proyecto}
\label{sec:anexo_repositorio}

El código fuente completo del proyecto está publicado en un repositorio
público de GitHub:

\begin{center}
  \url{https://github.com/hirojika/datawarehouse-transportes-rapiditos}
\end{center}

El repositorio reúne los scripts del Datawarehouse, las pruebas automáticas,
las herramientas de apoyo, la documentación del diseño y las fuentes de este
informe, de modo que cualquier persona puede revisar o reproducir el trabajo.
Cada entrega se registra en Git con un commit descriptivo, y el archivo
\bd{docs/bitacora.md} documenta, con fecha, cada decisión, cambio y
verificación realizada durante el proyecto. La Tabla~\ref{tab:repositorio}
describe su contenido.

\begin{table}[H]
  \centering
  \caption{Contenido del repositorio del proyecto}
  \label{tab:repositorio}
  \small
  \begin{tabularx}{\linewidth}{@{}>{\raggedright\arraybackslash}p{4.4cm}
                                  >{\raggedright\arraybackslash}X@{}}
    \toprule
    \textbf{Carpeta o archivo} & \textbf{Contenido} \\
    \midrule
    \bd{sql/} & Scripts del Datawarehouse: \bd{01_ddl.sql} (creación del
      modelo estrella) y \bd{99_validacion.sql} (verificación) \\
    \bd{docs/} & Diseño del modelo dimensional (\bd{modelo_dimensional.md},
      fuente de verdad del diseño), fuentes de información, bitácora,
      diagrama del esquema estrella y capturas del proceso \\
    \bd{informe/} & Fuentes LaTeX de este informe y de la guía de ejecución en
      Windows, con las referencias en \bd{referencias.bib} \\
    \bd{tests/} & Pruebas automáticas del script del Datawarehouse (pytest) \\
    \bd{herramientas/} & Captura automatizada de pgAdmin, generación del
      diccionario de datos (Anexo~A) y armado del paquete de entrega \\
    \bd{docker-compose.yml}, \bd{docker/} & Definición de los servicios
      PostgreSQL 17 y pgAdmin 4, con versiones fijas \\
    \bd{requirements.txt}, \bd{.env.example} & Dependencias de Python con
      versiones fijas y plantilla de configuración de la conexión \\
    \bd{etl/}, \bd{dashboard/}, \bd{data/} & Carpetas reservadas para el
      proyecto de Apache Hop, el dashboard en Streamlit y los datos de origen,
      que se desarrollan en los Avances N°2 y N°3 \\
    \bottomrule
  \end{tabularx}
\end{table}

Por seguridad y derechos de autor, el repositorio no incluye las credenciales
de conexión (archivo \bd{.env}, que cada integrante crea a partir de
\bd{.env.example}) ni el material del curso entregado por el profesor. Para
obtener una copia local se usa:

\begin{lstlisting}[style=consola]
git clone https://github.com/hirojika/datawarehouse-transportes-rapiditos.git
\end{lstlisting}

Los pasos para levantar el entorno a partir de esa copia están en el
Anexo~B (Sección~\ref{sec:anexo_entorno}).
